%% ============================================================
%% Teil III: rclpy – ROS2 mit Python (Kapitel 9–14)
%% Band 2 – Kapitel 9 bis 14
%% ============================================================

\part{rclpy -- ROS2 mit Python}

% ─────────────────────────────────────────────────────────────
\chapter{Wir wollen einen echten ROS2-Node schreiben}
% ─────────────────────────────────────────────────────────────

\begin{lernziele}
Du schreibst einen vollständigen ROS2-Node mit rclpy,
der sich korrekt initialisiert, loggt und wieder beendet.
Du verstehst den Unterschied zwischen \texttt{rclpy.spin()} und
\texttt{rclpy.spin\_once()} und weißt, was der Executor tut.
\end{lernziele}

\section{Situation}

In Band~1 Kapitel~27 hast du den ROS2-Node-Graph im Kopf aufgebaut.
Jetzt kommt der erste echte Node:
\texttt{arduino\_bridge\_node} bekommt seine ROS2-Grundstruktur.

\begin{aufgabeBox}
Schreibe den vollständigen \texttt{arduino\_bridge\_node.py}
als rclpy-Node, der startet, sich meldet und sauber beendet.
\end{aufgabeBox}

\section{Die rclpy-Grundstruktur}

\begin{wissenBox}
\begin{lstlisting}
#!/usr/bin/env python3
"""Minimaler rclpy-Node als Ausgangspunkt."""

import rclpy
from rclpy.node import Node

class MinimalNode(Node):
    def __init__(self) -> None:
        super().__init__('minimal_node')
        self.get_logger().info("MinimalNode gestartet.")
        # Timer: alle 1.0 Sekunden
        self._timer = self.create_timer(1.0, self._tick)
        self._zaehler: int = 0

    def _tick(self) -> None:
        self._zaehler += 1
        self.get_logger().info(f"Tick #{self._zaehler}")

def main(args: list[str] | None = None) -> None:
    rclpy.init(args=args)
    node = MinimalNode()
    try:
        rclpy.spin(node)      # Blockiert, bis Strg+C
    except KeyboardInterrupt:
        pass
    finally:
        node.destroy_node()   # ROS2-Ressourcen freigeben
        rclpy.shutdown()      # rclpy herunterfahren

if __name__ == '__main__':
    main()
\end{lstlisting}

\textbf{\texttt{rclpy.spin(node)}:} Blockiert und verarbeitet alle
eingehenden Nachrichten und Timer-Ticks -- bis der Node zerstört wird.

\textbf{\texttt{rclpy.spin\_once(node)}:} Verarbeitet genau eine Nachricht
oder einen Timer-Tick, dann kehrt die Funktion zurück.
Nützlich, wenn du die Kontrolle zwischendurch brauchst.
\end{wissenBox}

\subsection{Was ist ein Executor?}

\begin{haubeBox}
\texttt{rclpy.spin(node)} erstellt intern einen \texttt{SingleThreadedExecutor}:
Er verarbeitet alle Callbacks (Timer, Subscriber, Service-Antworten)
nacheinander in einem einzigen Thread.

Für fortgeschrittene Nodes mit vielen parallelen Callbacks
gibt es \texttt{MultiThreadedExecutor}:
\begin{lstlisting}
from rclpy.executors import MultiThreadedExecutor
executor = MultiThreadedExecutor(num_threads=4)
executor.add_node(node)
executor.spin()
\end{lstlisting}
Für \texttt{arduino\_bridge\_node} reicht \texttt{SingleThreadedExecutor} --
Sensor lesen und publishen ist sequenziell.
\end{haubeBox}

\section{Schritt-für-Schritt-Lösung}

\begin{loesungBox}
\begin{lstlisting}[language=bash]
# Testen ohne ROS2-Paket (direkt)
cd ~/robotik-uno-q/stufe2_ros2
source /opt/ros/jazzy/setup.bash
python3 src/arduino_bridge_node.py

# In einem zweiten Terminal:
source /opt/ros/jazzy/setup.bash
ros2 node list
# /arduino_bridge
\end{lstlisting}
\end{loesungBox}

\begin{projektBox}[robotik-uno-q -- Projektstand]
\textbf{Heute:} Erster echter ROS2-Node läuft. \texttt{ros2 node list} zeigt ihn.\\
\textbf{Nächstes Kapitel:} Sensordaten als Topic publishen.
\end{projektBox}

\begin{merksatzBox}
\textbf{Merksatz:}
\texttt{rclpy.init()} → Node erstellen → \texttt{rclpy.spin()} → \texttt{destroy\_node()} → \texttt{rclpy.shutdown()}.
Diese fünf Schritte sind das Grundgerüst jedes rclpy-Nodes.
\texttt{try/except KeyboardInterrupt} fängt Strg+C sauber ab.
\end{merksatzBox}

\begin{fehlerBox}
\begin{itemize}
\item \texttt{rclpy.init()} zweimal aufrufen: \texttt{RuntimeError: Already initialized}.
\item \texttt{destroy\_node()} vergessen: Port-Konflikte beim nächsten Start.
\item Node-Name mit Bindestrich: \texttt{'arduino-bridge'} → Fehler.
  Nur Buchstaben, Ziffern und Unterstriche erlaubt.
\end{itemize}
\end{fehlerBox}


% ─────────────────────────────────────────────────────────────
\chapter{Wir wollen Sensordaten publishen und empfangen –
  Publisher und Subscriber}
% ─────────────────────────────────────────────────────────────

\begin{lernziele}
Du implementierst \texttt{arduino\_bridge\_node} als Publisher
für \texttt{sensor\_msgs/Range} und \texttt{navigation\_node}
als Subscriber -- und verstehst QoS-Profile.
\end{lernziele}

\section{Situation}

\texttt{arduino\_bridge\_node} liest Sensordaten.
\texttt{navigation\_node} braucht sie.
In Band~1 haben wir dafür den \texttt{SensorBus} gebaut.
ROS2-Topics ersetzen ihn -- mit Netzwerkfähigkeit, Zeitstempeln
und Standard-Nachrichtentypen.

\begin{aufgabeBox}
Implementiere Publisher für \texttt{/sensor/rechts} und \texttt{/sensor/links}
(\texttt{sensor\_msgs/Range}) im \texttt{arduino\_bridge\_node}
und einen Subscriber im \texttt{navigation\_node}.
\end{aufgabeBox}

\section{Nachrichtentypen: sensor\_msgs/Range}

\begin{lstlisting}
from sensor_msgs.msg import Range

# Range-Nachricht füllen
msg = Range()
msg.header.stamp = self.get_clock().now().to_msg()
msg.header.frame_id = 'sensor_rechts'
msg.radiation_type  = Range.ULTRASOUND    # HC-SR04
msg.field_of_view   = 0.26               # ~15 Grad in Radiant
msg.min_range       = 0.02               # 2 cm
msg.max_range       = 4.00               # 400 cm
msg.range           = 0.87              # aktueller Messwert in Metern
\end{lstlisting}

\section{Publisher implementieren}

\begin{lstlisting}
from rclpy.node import Node
from rclpy.qos import QoSProfile, ReliabilityPolicy, HistoryPolicy
from sensor_msgs.msg import Range

class ArduinoBridgeNode(Node):
    def __init__(self) -> None:
        super().__init__('arduino_bridge')

        # QoS: Sensor-Daten -- Best Effort, nur letzter Wert zählt
        sensor_qos = QoSProfile(
            reliability=ReliabilityPolicy.BEST_EFFORT,
            history=HistoryPolicy.KEEP_LAST,
            depth=1
        )

        self._pub_rechts = self.create_publisher(
            Range, '/sensor/rechts', sensor_qos)
        self._pub_links  = self.create_publisher(
            Range, '/sensor/links',  sensor_qos)

        self._timer = self.create_timer(0.1, self._lese_und_publiziere)
        self.get_logger().info("ArduinoBridgeNode bereit.")

    def _lese_und_publiziere(self) -> None:
        # Aus Band 1: Serial lesen und parsen
        zeile = self._lese_serielle_zeile()
        if zeile is None:
            return
        rechts_cm, links_cm = zeile

        jetzt = self.get_clock().now().to_msg()
        self._pub_rechts.publish(self._baue_range_msg(rechts_cm, 'sensor_rechts', jetzt))
        self._pub_links.publish( self._baue_range_msg(links_cm,  'sensor_links',  jetzt))

    def _baue_range_msg(self, cm: int, frame: str, stamp) -> Range:
        msg = Range()
        msg.header.stamp     = stamp
        msg.header.frame_id  = frame
        msg.radiation_type   = Range.ULTRASOUND
        msg.field_of_view    = 0.26
        msg.min_range        = 0.02
        msg.max_range        = 4.00
        msg.range            = float(cm) / 100.0
        return msg
\end{lstlisting}

\section{Subscriber implementieren}

\begin{lstlisting}
from rclpy.node import Node
from sensor_msgs.msg import Range

class NavigationNode(Node):
    def __init__(self) -> None:
        super().__init__('navigation')

        sensor_qos = QoSProfile(
            reliability=ReliabilityPolicy.BEST_EFFORT,
            history=HistoryPolicy.KEEP_LAST,
            depth=1
        )
        self._sub_rechts = self.create_subscription(
            Range, '/sensor/rechts', self._cb_rechts, sensor_qos)
        self._sub_links  = self.create_subscription(
            Range, '/sensor/links',  self._cb_links,  sensor_qos)

        self._rechts_m: float = 4.0
        self._links_m:  float = 4.0

    def _cb_rechts(self, msg: Range) -> None:
        self._rechts_m = msg.range
        self._entscheide()

    def _cb_links(self, msg: Range) -> None:
        self._links_m = msg.range
        self._entscheide()

    def _entscheide(self) -> None:
        minimum = min(self._rechts_m, self._links_m)
        if minimum < 0.25:
            self.get_logger().warning(f"STOPP: {minimum:.2f}m")
        elif minimum < 0.50:
            self.get_logger().info(f"LANGSAM: {minimum:.2f}m")
        else:
            self.get_logger().debug(f"FREI: R={self._rechts_m:.2f}m L={self._links_m:.2f}m")
\end{lstlisting}

\begin{haubeBox}
\textbf{QoS -- Quality of Service}

\texttt{RELIABLE}: Jede Nachricht wird garantiert zugestellt -- wie TCP.
Nötig für Service-Aufrufe und kritische Befehle.

\texttt{BEST\_EFFORT}: Pakete können verloren gehen -- wie UDP.
Ideal für Sensordaten: ein verlorener Messwert ist unwichtig,
Latenz dagegen kritisch.

\texttt{KEEP\_LAST + depth=1}: Subscriber empfängt nur den neuesten Wert.
Für Sensordaten fast immer richtig: alte Abstandswerte sind wertlos.
\end{haubeBox}

\begin{projektBox}[robotik-uno-q -- Projektstand]
\textbf{Heute:} \texttt{arduino\_bridge\_node} published \texttt{sensor\_msgs/Range}.
\texttt{navigation\_node} subscribed und trifft erste Fahrentscheidungen.\\
\textbf{Nächstes Kapitel:} Motor-Befehle über \texttt{/cmd\_vel} senden.
\end{projektBox}

\begin{merksatzBox}
\textbf{Merksatz:}
\texttt{BEST\_EFFORT} für Sensor-Streams.
\texttt{RELIABLE} für Service-Calls und Motor-Befehle.
Publisher und Subscriber müssen \emph{kompatible} QoS-Profile haben,
sonst fließen keine Nachrichten.
\end{merksatzBox}


% ─────────────────────────────────────────────────────────────
\chapter{Wir wollen Motor-Befehle senden – \texttt{geometry\_msgs/Twist}}
% ─────────────────────────────────────────────────────────────

\begin{lernziele}
Du implementierst den \texttt{/cmd\_vel}-Publisher im \texttt{navigation\_node}
mit \texttt{geometry\_msgs/Twist}, verstehst das Kinematikmodell
des Differentialantriebs und weißt,
wie \texttt{motor\_node} den Twist-Befehl umsetzt.
\end{lernziele}

\section{Situation}

\texttt{navigation\_node} hat entschieden: links ausweichen.
Jetzt braucht er einen Weg, dem \texttt{motor\_node}
"`fahre links"' zu sagen -- in ROS2-Standardform.

\begin{aufgabeBox}
Erweitere \texttt{navigation\_node} um einen \texttt{/cmd\_vel}-Publisher
und implementiere die Entscheidungslogik für das Ausweichmanöver.
\end{aufgabeBox}

\section{geometry\_msgs/Twist}

\begin{wissenBox}
\texttt{Twist} beschreibt eine Geschwindigkeit im 3D-Raum.
Für einen mobilen Roboter auf dem Boden sind nur zwei Felder relevant:

\begin{lstlisting}
from geometry_msgs.msg import Twist

msg = Twist()
msg.linear.x  = 0.3   # Vorwärts in m/s  (>0 vorwärts, <0 rückwärts)
msg.angular.z = 0.0   # Rotation in rad/s (>0 links, <0 rechts)

# Beispiele:
# Geradeaus 0.2 m/s:  linear.x=0.2, angular.z=0.0
# Linkskurve:         linear.x=0.1, angular.z=0.5
# Auf der Stelle:     linear.x=0.0, angular.z=1.0
# Stopp:              linear.x=0.0, angular.z=0.0
\end{lstlisting}
\end{wissenBox}

\begin{lstlisting}
class NavigationNode(Node):
    # ... (Subscriber aus Kapitel 10 behalten)

    def __init__(self) -> None:
        super().__init__('navigation')
        # ... Subscriber ...
        self._pub_cmd = self.create_publisher(
            Twist, '/cmd_vel', 10)

    def _entscheide(self) -> None:
        cmd = Twist()
        minimum = min(self._rechts_m, self._links_m)

        if minimum < 0.25:                   # Stopp
            cmd.linear.x  = 0.0
            cmd.angular.z = 0.0
        elif minimum < 0.50:                 # Langsam + ausweichen
            cmd.linear.x  = 0.1
            cmd.angular.z = 0.8 if self._rechts_m < self._links_m else -0.8
        else:                                # Geradeaus
            cmd.linear.x  = 0.2
            cmd.angular.z = 0.0

        self._pub_cmd.publish(cmd)
\end{lstlisting}

\begin{projektBox}[robotik-uno-q -- Projektstand]
\textbf{Heute:} \texttt{navigation\_node} published \texttt{/cmd\_vel}.\\
\textbf{Nächstes Kapitel:} Services -- Notbremse als \texttt{/motor/stopp}.
\end{projektBox}

\begin{merksatzBox}
\textbf{Merksatz:} \texttt{linear.x} = Translation (vorwärts/rückwärts),
\texttt{angular.z} = Rotation (links/rechts).
Alle anderen Felder bleiben 0.0 für einen planaren Roboter.
Twist ist der universelle ROS2-Befehl für mobile Roboter --
Nav2 sendet ihn, \texttt{ros2 teleop\_twist\_keyboard} sendet ihn,
dein Navigation-Node sendet ihn.
\end{merksatzBox}


% ─────────────────────────────────────────────────────────────
\chapter{Wir wollen eine Notbremse – Services in rclpy}
% ─────────────────────────────────────────────────────────────

\begin{lernziele}
Du implementierst \texttt{/motor/stopp} als ROS2-Service,
verstehst den Unterschied zwischen Topic und Service
und weißt, wann welches Pattern passt.
\end{lernziele}

\section{Situation}

Topics sind asynchron: Der Publisher sendet, ohne zu wissen,
ob jemand empfängt. Für eine Notbremse ist das zu unsicher.
Services garantieren eine Antwort: "`Stopp-Befehl empfangen"' oder
"`Fehler: Motor nicht erreichbar"'.

\begin{aufgabeBox}
Implementiere \texttt{/motor/stopp} als \texttt{std\_srvs/Trigger}-Service
im \texttt{motor\_node} und rufe ihn aus dem \texttt{navigation\_node} auf.
\end{aufgabeBox}

\section{Service-Server implementieren}

\begin{lstlisting}
from std_srvs.srv import Trigger
from rclpy.node import Node

class MotorNode(Node):
    def __init__(self) -> None:
        super().__init__('motor')
        self._srv = self.create_service(
            Trigger, '/motor/stopp', self._stopp_callback)
        self.get_logger().info("MotorNode bereit. Service /motor/stopp aktiv.")

    def _stopp_callback(
        self,
        request:  Trigger.Request,
        response: Trigger.Response
    ) -> Trigger.Response:
        """Notbremse: sendet Stopp-Befehl an Arduino."""
        try:
            # Hier: 'S\n' über Serial senden (aus Band 1)
            self._sende_arduino('S')
            response.success = True
            response.message = "Motoren gestoppt."
            self.get_logger().warning("NOTBREMSE ausgeloest.")
        except Exception as e:
            response.success = False
            response.message = str(e)
        return response
\end{lstlisting}

\section{Service-Client aufrufen}

\begin{lstlisting}
from std_srvs.srv import Trigger

class NavigationNode(Node):
    def __init__(self) -> None:
        # ...
        self._stopp_client = self.create_client(Trigger, '/motor/stopp')

    def _notbremse(self) -> None:
        """Synchroner Stopp-Aufruf (blockiert kurz)."""
        if not self._stopp_client.wait_for_service(timeout_sec=1.0):
            self.get_logger().error("Service /motor/stopp nicht erreichbar!")
            return
        req = Trigger.Request()
        future = self._stopp_client.call_async(req)
        # In einem einfachen Node: auf Antwort warten
        rclpy.spin_until_future_complete(self, future, timeout_sec=2.0)
        if future.result().success:
            self.get_logger().warning("Notbremse bestaetigt.")
\end{lstlisting}

\begin{projektBox}[robotik-uno-q -- Projektstand]
\textbf{Heute:} \texttt{/motor/stopp}-Service implementiert und getestet.\\
\textbf{Nächstes Kapitel:} Parameter -- Schwellenwerte zur Laufzeit ändern.
\end{projektBox}

\begin{merksatzBox}
\textbf{Merksatz:}
Topic = "`Ich schicke etwas, egal ob jemand zuhört."'
Service = "`Ich schicke etwas und warte auf Bestätigung."'
Notbremsen, Konfigurationsänderungen und einmalige Aktionen:
Service. Sensordaten, Befehle mit hoher Frequenz: Topic.
\end{merksatzBox}


% ─────────────────────────────────────────────────────────────
\chapter{Wir wollen Parameter zur Laufzeit ändern –
  \texttt{declare\_parameter} und Konfiguration}
% ─────────────────────────────────────────────────────────────

\begin{lernziele}
Du verwendest ROS2-Parameter statt hartcodierter Werte,
lädst \texttt{sensor\_config.yaml} über die Parameter-API
und kannst Werte zur Laufzeit mit \texttt{ros2 param set} ändern.
\end{lernziele}

\section{Situation}

\texttt{sensor\_bridge.py} aus Band~1 lädt \texttt{sensor\_config.yaml}
direkt mit Python's yaml-Modul.
In ROS2 gibt es einen besseren Weg:
Parameter werden beim Node-Start übergeben,
können zur Laufzeit geändert werden
und erscheinen in \texttt{ros2 param list}.

\begin{aufgabeBox}
Ersetze die YAML-Direktladung durch ROS2-Parameter in \texttt{arduino\_bridge\_node}.
\end{aufgabeBox}

\begin{lstlisting}
class ArduinoBridgeNode(Node):
    def __init__(self) -> None:
        super().__init__('arduino_bridge')

        # Parameter deklarieren (mit Defaultwert)
        self.declare_parameter('port',          '/dev/arduino')
        self.declare_parameter('baudrate',      9600)
        self.declare_parameter('min_abstand_m', 0.25)
        self.declare_parameter('warn_abstand_m',0.50)

        # Parameter lesen
        self._port     = self.get_parameter('port').value
        self._baudrate = self.get_parameter('baudrate').value
        self._min_m    = self.get_parameter('min_abstand_m').value

        self.get_logger().info(
            f"Parameter: port={self._port}, baudrate={self._baudrate}")
\end{lstlisting}

\begin{lstlisting}[language=bash]
# Node mit Parameters-File starten
ros2 run robotik_uno_q arduino_bridge_node \
    --ros-args --params-file config/sensor_config.yaml

# Zur Laufzeit ändern
ros2 param set /arduino_bridge min_abstand_m 0.30

# Parameter anzeigen
ros2 param list /arduino_bridge
ros2 param get /arduino_bridge baudrate
\end{lstlisting}

\begin{projektBox}[robotik-uno-q -- Projektstand]
\textbf{Heute:} Parameter statt YAML-Direktladung. Schwellenwerte zur Laufzeit änderbar.\\
\textbf{Nächstes Kapitel:} Launch-Datei -- alle Nodes mit einem Befehl starten.
\end{projektBox}

\begin{merksatzBox}
\textbf{Merksatz:}
\texttt{declare\_parameter('name', default)} registriert einen Parameter mit Fallback.
\texttt{--params-file sensor\_config.yaml} überschreibt Defaults beim Start.
\texttt{ros2 param set} ändert Werte zur Laufzeit -- ohne Neustart.
\end{merksatzBox}


% ─────────────────────────────────────────────────────────────
\chapter{Wir wollen alle Nodes gemeinsam starten – Launch-Dateien}
% ─────────────────────────────────────────────────────────────

\begin{lernziele}
Du schreibst eine Python-Launch-Datei, die alle drei Nodes
(\texttt{arduino\_bridge}, \texttt{navigation}, \texttt{motor}) mit
korrekten Parametern startet -- ein Befehl für den kompletten Stack.
\end{lernziele}

\section{Situation}

Drei Terminal-Fenster für drei Nodes -- das ist kein Deployment.
\texttt{start\_roboter.sh} aus Band~1 hat dieses Problem gelöst.
ROS2 löst es sauberer: mit Launch-Dateien,
die Parameter, Namespaces und Log-Level kennen.

\begin{aufgabeBox}
Schreibe \texttt{stufe2\_ros2/launch/roboter.launch.py},
das alle drei Nodes startet und \texttt{sensor\_config.yaml} übergibt.
\end{aufgabeBox}

\begin{lstlisting}
# stufe2_ros2/launch/roboter.launch.py
"""Launch-Datei für den vollständigen robotik-uno-q Stack."""

from launch import LaunchDescription
from launch_ros.actions import Node
from launch.actions import DeclareLaunchArgument
from launch.substitutions import LaunchConfiguration
from pathlib import Path

PAKET = 'robotik_uno_q'
CONFIG = Path(__file__).parent.parent / 'config' / 'sensor_config.yaml'

def generate_launch_description() -> LaunchDescription:
    port_arg = DeclareLaunchArgument(
        'port',
        default_value='/dev/arduino',
        description='Serial-Port des Arduino'
    )

    arduino_bridge = Node(
        package=PAKET,
        executable='arduino_bridge_node',
        name='arduino_bridge',
        parameters=[str(CONFIG), {'port': LaunchConfiguration('port')}],
        output='screen',
    )

    navigation = Node(
        package=PAKET,
        executable='navigation_node',
        name='navigation',
        parameters=[str(CONFIG)],
        output='screen',
    )

    motor = Node(
        package=PAKET,
        executable='motor_node',
        name='motor',
        parameters=[str(CONFIG), {'port': LaunchConfiguration('port')}],
        output='screen',
    )

    return LaunchDescription([
        port_arg,
        arduino_bridge,
        navigation,
        motor,
    ])
\end{lstlisting}

\begin{lstlisting}[language=bash]
# Stack starten
ros2 launch robotik_uno_q roboter.launch.py

# Mit anderem Port:
ros2 launch robotik_uno_q roboter.launch.py port:=/dev/ttyACM0

# Node-Graph anzeigen
rqt_graph
\end{lstlisting}

\begin{projektBox}[robotik-uno-q -- Projektstand]
\textbf{Heute:} \texttt{roboter.launch.py} -- ein Befehl startet den kompletten Stack.\\
\textbf{Teil~IV:} Alles zusammensetzen -- vollständige Implementierung.
\end{projektBox}

\begin{merksatzBox}
\textbf{Merksatz:}
\texttt{generate\_launch\_description()} gibt die Liste aller zu startenden Actions zurück.
\texttt{DeclareLaunchArgument} + \texttt{LaunchConfiguration} machen Parameter
über die Kommandozeile überschreibbar.
Launch-Dateien ersetzen \texttt{start\_roboter.sh} -- aber das Shell-Skript
bleibt für Docker-Umgebungen nützlich.
\end{merksatzBox}

\begin{haubeBox}
\textbf{Warum Python statt XML für Launch-Dateien?}

ROS1 verwendete XML-Launch-Dateien (\texttt{.launch}).
Sie waren einfach, aber unflexibel: keine Schleifen,
keine bedingten Ausdrücke, keine Imports.
ROS2 führte Python-Launch-Dateien ein -- volle Programmierbarkeit,
Zugriff auf das Filesystem, bedingte Nodes, Schleifen über Node-Listen.
XML-Launch-Dateien gibt es in ROS2 noch, aber die Python-API ist Standard.
\end{haubeBox}
